% ---------------------------------------------------------------- Cryptographic implementation
\section{Cryptographic Implementation}
\label{sec:crypto}

All cryptography is written from scratch in the \texttt{rsh-crypto} crate,
about 410 lines with no dependencies. This section explains each part and then
follows one real recorded session byte by byte. The worked values come from a
session between the real \texttt{rsh} server and client, recorded through the
TCP relay of Section~\ref{sec:telnetwire} and decoded afterwards with the
server's private key (\texttt{report/tools/decoder}). Every table in
Sections~\ref{sec:hostkey}--\ref{sec:packets} is generated from that recording
(\texttt{report/data/}), and every value was checked by the decoder against the
library code itself.

\begin{table}[H]
\centering
\footnotesize
\begin{tabularx}{\textwidth}{@{}l l X@{}}
\toprule
Building block & File & Used for \\
\midrule
Modular arithmetic & \texttt{math.rs} & $a b \bmod m$, $a^{k} \bmod m$, $a^{-1} \bmod m$ on 64-bit numbers \\
Random numbers & \texttt{random.rs} & candidate primes, the 32-byte session key \\
Miller--Rabin & \texttt{prime.rs} & testing candidate primes \\
RSA, 64-bit modulus & \texttt{rsa.rs} & sending the session key; signing the login \\
Digest $H$ (FNV-1a) & \texttt{hash.rs} & the value that is signed; key fingerprints \\
Stream cipher & \texttt{cipher.rs} & encrypting every message after the key exchange \\
Key files & \texttt{keyfile.rs} & \texttt{id\_rsh}, \texttt{host\_key}, \texttt{authorized\_keys}, \texttt{known\_hosts} \\
\bottomrule
\end{tabularx}
\caption{The cryptographic building blocks.}
\end{table}

RSA is slow and works only on numbers smaller than $n$, so it is used twice per
connection: once to send the session key and once to sign the login. All the
terminal traffic is then encrypted with the fast stream cipher. Real SSH divides
the work the same way: public-key operations only to agree on a key and to
authenticate, a symmetric cipher for the data.

% ................................................................
\subsection{Modular arithmetic}

Every RSA operation is arithmetic modulo $n < 2^{64}$. The product of two such
numbers needs up to 128 bits, so \texttt{mul\_mod} widens to \texttt{u128}
before reducing:
\[
  \texttt{mul\_mod}(a, b, m) = (a \cdot b) \bmod m, \qquad a, b < m < 2^{64}, \quad a \cdot b < 2^{128}.
\]
\texttt{pow\_mod} computes $b^{k} \bmod m$ by \emph{square and multiply}. It
reads the exponent bit by bit from the lowest bit up: the base is squared at
every step and multiplied into the result where the bit is 1. A 64-bit exponent
therefore costs at most 64 squarings and 64 multiplications, instead of $k$
multiplications. \texttt{inv\_mod} finds $a^{-1} \bmod m$ with the extended
Euclidean algorithm, which keeps $s$ such that $\mathit{old\_r} \equiv
\mathit{old\_s} \cdot a \pmod m$ at every step.

\Needspace{24\baselineskip}
\begin{lstlisting}[caption={Square and multiply, and the modular inverse (\texttt{math.rs}).}]
pub fn pow_mod(mut base: u64, mut exp: u64, m: u64) -> u64 {
    if m == 1 { return 0; }
    let mut result = 1;
    base %= m;
    while exp > 0 {
        if exp & 1 == 1 { result = mul_mod(result, base, m); }   // bit is 1: multiply
        base = mul_mod(base, base, m);                            // square for the next bit
        exp >>= 1;
    }
    result
}
pub fn inv_mod(a: u64, m: u64) -> Option<u64> {
    let (mut old_r, mut r) = (a as i128, m as i128);
    let (mut old_s, mut s) = (1i128, 0i128);
    while r != 0 {
        let q = old_r / r;
        (old_r, r) = (r, old_r - q * r);
        (old_s, s) = (s, old_s - q * s);
    }
    if old_r != 1 { return None; }                  // gcd(a, m) != 1: no inverse
    Some(old_s.rem_euclid(m as i128) as u64)
}
\end{lstlisting}

% ................................................................
\subsection{RSA by hand, with a small key}

The same functions work on small numbers, which makes them easy to check by
hand. Take $p = 61$ and $q = 53$:
\[
  n = 61 \cdot 53 = 3233, \qquad \varphi = 60 \cdot 52 = 3120, \qquad e = 65537 .
\]
\textbf{Private exponent.} \texttt{inv\_mod(65537, 3120)} runs as in
Table~\ref{tab:toyeuclid} and gives $d = -367 \bmod 3120 = 2753$. Check:
$65537 \cdot 2753 = 180\,423\,361 = 57\,828 \cdot 3120 + 1$.

\begin{table}[H]
\centering
\small
\begin{tabular}{@{}r r r r r@{}}
\toprule
$\mathit{old\_r}$ & $r$ & $q = \lfloor \mathit{old\_r} / r \rfloor$ & $\mathit{old\_s}$ & $s$ \\
\midrule
65537 & 3120 & 21  & 1    & 0 \\
3120  & 17   & 183 & 0    & 1 \\
17    & 9    & 1   & 1    & $-183$ \\
9     & 8    & 1   & $-183$ & 184 \\
8     & 1    & 8   & 184  & $-367$ \\
1     & 0    & --  & $-367$ & 3120 \\
\bottomrule
\end{tabular}
\caption{Extended Euclid for $65537^{-1} \bmod 3120$. It stops with $\mathit{old\_r} = 1 = \gcd$, so the inverse exists.}
\label{tab:toyeuclid}
\end{table}

\textbf{Encryption.} Encrypt the letter \texttt{A}, $m = 65$. The exponent
$65537 = 2^{16} + 1$ has only two 1-bits, so \texttt{pow\_mod} squares 16 times
and multiplies twice: $c = 65^{2^{16}} \cdot 65 \bmod 3233$.

\begin{table}[H]
\centering
\small
\setlength{\tabcolsep}{3.2pt}
\begin{tabular}{@{}l *{17}{r}@{}}
\toprule
$k$ & 0 & 1 & 2 & 3 & 4 & 5 & 6 & 7 & 8 & 9 & 10 & 11 & 12 & 13 & 14 & 15 & 16 \\
$65^{2^k} \bmod 3233$ & 65 & 992 & 1232 & 1547 & 789 & 1785 & 1720 & 205 & 3229 & 16 & 256 & 876 & 1155 & 2029 & 1232 & 1547 & 789 \\
bit $k$ of $e$ & 1 & 0 & 0 & 0 & 0 & 0 & 0 & 0 & 0 & 0 & 0 & 0 & 0 & 0 & 0 & 0 & 1 \\
\bottomrule
\end{tabular}
\caption{Repeated squaring of $65$ modulo $3233$.}
\end{table}

\[
  c = 789 \cdot 65 \bmod 3233 = 51\,285 \bmod 3233 = 2790, \qquad
  m = 2790^{2753} \bmod 3233 = 65 = \texttt{A}.
\]
Decryption uses the same \texttt{pow\_mod} with $d = 2753 =
\texttt{101011000001}_2$ (12 squarings, 5 multiplications) and gives back 65.
This works because $e d \equiv 1 \pmod{\varphi}$, so
$m^{ed} = m^{1 + k\varphi} \equiv m \pmod n$ by Euler's theorem.

% ................................................................
\subsection{Making a key pair}
\label{sec:hostkey}

\texttt{KeyPair::generate} does the same with two random 32-bit primes.
\texttt{random\_prime} takes a 64-bit random number, keeps the low 32 bits,
sets the top bit (so that the prime really has 32 bits) and the low bit (so
that it is odd), and tests it with Miller--Rabin until one is prime. Setting the
top bit of both primes gives $n \geq 2^{62}$, so every 7-byte block
($< 2^{56}$) is smaller than $n$. The host key of the recorded session:

\begin{table}[H]
\centering
\small
% generated by tools/gen_tables.py from data/ -- do not edit
\begin{tabular}{@{}l r l@{}}
\toprule
Quantity & Value & How it was obtained \\
\midrule
$p$ & 3\,704\,941\,837 & random 32-bit odd number, top bit set; prime on try 1 \\
$q$ & 4\,252\,763\,287 & the same; prime on try 2 \\
$n = pq$ & 15\,756\,240\,624\,863\,938\,219 & 64 bits, \texttt{daa9695e409806ab} in hex \\
$\varphi = (p-1)(q-1)$ & 15\,756\,240\,616\,906\,233\,096 & \\
$e$ & 65537 & fixed, $\gcd(e, \varphi) = 1$ \\
$d = e^{-1} \bmod \varphi$ & 12\,084\,825\,226\,199\,074\,001 & extended Euclid; check $e d \bmod \varphi = 1$ \\
\bottomrule
\end{tabular}

\caption{The host key of the recorded session, made by the same steps as \texttt{KeyPair::generate}.}
\label{tab:hostkey}
\end{table}

\textbf{Miller--Rabin.} Write $p - 1 = 2^{s} d$ with $d$ odd. If $p$ is prime,
then for every witness $a$ the sequence $a^{d}, a^{2d}, \ldots, a^{2^{s-1}d}
\pmod p$ either starts at 1 or reaches $p - 1$. A composite number fails this
for most $a$, and with the first twelve primes $2, 3, \ldots, 37$ as witnesses no composite
below $3.1 \cdot 10^{23}$ passes, so the test is exact for every 64-bit number. For $p$ above:

\begin{table}[H]
\centering
\small
% generated by tools/gen_tables.py from data/ -- do not edit
\begin{tabularx}{\textwidth}{@{}c X l@{}}
\toprule
$a$ & Squarings, with $p - 1 = 2^{2} \cdot 926\,235\,459$ & Result \\
\midrule
2 & $x_0 = 2^{d} \bmod p = 2\,338\,420\,328$; $x_1 = x_0^2 \bmod p = 3\,704\,941\,836$ & $= p-1$, passes \\
3 & $x_0 = 3^{d} \bmod p = 3\,704\,941\,836$ & $= p-1$, passes \\
5 & $x_0 = 5^{d} \bmod p = 1\,366\,521\,509$; $x_1 = x_0^2 \bmod p = 3\,704\,941\,836$ & $= p-1$, passes \\
\bottomrule
\end{tabularx}

\caption{The first three of the twelve Miller--Rabin witnesses for $p$.}
\end{table}

% ................................................................
\subsection{Encrypting the session key with RSA}
\label{sec:rsablocks}

The client makes a random 32-byte session key $K$. In the recorded session:
\begin{center}
% generated by tools/gen_tables.py from data/ -- do not edit
\texttt{53 51 46 f5 60 34 ab 63 20 13 66 85 7b 35 96 3e}\\
\texttt{df f5 68 45 22 dd 29 25 31 81 e6 d2 53 79 7d d0}
\end{center}
\texttt{PublicKey::encrypt} writes the plaintext length (4 bytes), then cuts the
data into 7-byte blocks. Each block is placed in the low 7 bytes of a 64-bit
number $m_i$, whose top byte stays 0, so that $m_i < 2^{56} \leq n$. Each block
is encrypted separately, $c_i = m_i^{e} \bmod n$, and written as 8 bytes. The
last block is padded with zeros; the length field tells the receiver how many
bytes to keep. 32 key bytes therefore need $\lceil 32/7 \rceil = 5$ blocks, or
$4 + 5 \cdot 8 = 44$ bytes.

\begin{table}[H]
\centering
\footnotesize
% generated by tools/gen_tables.py from data/ -- do not edit
\begin{tabular}{@{}c l l l c@{}}
\toprule
$i$ & Key bytes in the block & $m_i$ (00 + 7 bytes) & $c_i = m_i^{e} \bmod n$ (on the wire) & $c_i^{d} \bmod n = m_i$ \\
\midrule
0 & \texttt{53 51 46 f5 60 34 ab} & \texttt{00535146f56034ab} & \texttt{239588b8167bdb4e} & yes \\
1 & \texttt{63 20 13 66 85 7b 35} & \texttt{0063201366857b35} & \texttt{825971f61e748762} & yes \\
2 & \texttt{96 3e df f5 68 45 22} & \texttt{00963edff5684522} & \texttt{52e5fdb92a92275d} & yes \\
3 & \texttt{dd 29 25 31 81 e6 d2} & \texttt{00dd29253181e6d2} & \texttt{7a0d14922e2c96dd} & yes \\
4 & \texttt{53 79 7d d0} & \texttt{0053797dd0000000} & \texttt{3a57496cec72280b} & yes \\
\bottomrule
\end{tabular}

\caption{The session key under RSA with the host key of Table~\ref{tab:hostkey}.}
\label{tab:rsablocks}
\end{table}

The \texttt{SessionKey} message carries the 44 bytes as a byte field, so the
whole frame is 53 bytes. It is the second frame of the connection and is not
encrypted by the stream cipher, which does not exist yet:

% generated by tools/gen_tables.py from data/ -- do not edit
\begin{lstlisting}[style=plain]
frame length   00 00 00 31                      = 49
tag            02                               = SessionKey
field length   00 00 00 2c                      = 44
plain length   00 00 00 20                      = 32 key bytes
c_0            23 95 88 b8 16 7b db 4e
c_1            82 59 71 f6 1e 74 87 62
c_2            52 e5 fd b9 2a 92 27 5d
c_3            7a 0d 14 92 2e 2c 96 dd
c_4            3a 57 49 6c ec 72 28 0b
\end{lstlisting}


The server runs \texttt{PrivateKey::decrypt}: for each $c_i$ it computes
$c_i^{d} \bmod n$ and keeps bytes 1--7, and from the last block only as many
bytes as the length field still needs. It ends up with the same 32 bytes.
An observer sees only $c_0, \ldots, c_4$ and, without $d$, cannot undo the
exponentiation (Section~\ref{sec:cryptolimits} shows that with a 64-bit $n$ this
is not true in practice).

\Needspace{14\baselineskip}
\begin{lstlisting}[caption={Block encryption (\texttt{rsa.rs}).}]
pub fn encrypt(&self, data: &[u8]) -> Vec<u8> {
    let mut out = (data.len() as u32).to_be_bytes().to_vec();   // plaintext length
    for chunk in data.chunks(PLAIN_BLOCK) {                      // 7 bytes at a time
        let mut block = [0u8; 8];
        block[1..1 + chunk.len()].copy_from_slice(chunk);        // top byte stays 0
        let m = u64::from_be_bytes(block);
        out.extend_from_slice(&pow_mod(m, self.e, self.n).to_be_bytes());
    }
    out
}
\end{lstlisting}

% ................................................................
\subsection{The digest and the signature}
\label{sec:signature}

\textbf{Digest.} Signatures are made over a 64-bit digest $H(x)$ instead of
the data itself. It is built from FNV-1a, which processes one byte at a time:
\[
  h_0 = \texttt{cbf29ce484222325}_{16}, \qquad
  h_{i+1} = (h_i \oplus x_i) \cdot \texttt{100000001b3}_{16} \bmod 2^{64}.
\]
To mix the bits further, the digest hashes the data a second time with the
first result appended and combines the two:
\[
  f = \mathrm{FNV1a}(x), \qquad
  H(x) = \mathrm{rotl}_{17}\bigl(\mathrm{FNV1a}(x \,\Vert\, f)\bigr) \oplus f .
\]
The fingerprint of a key, as printed by \texttt{rsh-keygen} and on first
connection, is $H(n \Vert e)$ written as eight hex bytes, for example
\texttt{0d:28:4c:ec:4e:a1:d9:04} for the host key above.

\textbf{Signature.} RSA signing is decryption with the private key applied to
the digest: $\sigma = (H(x) \bmod n)^{d} \bmod n$. Anyone with the public key
checks $\sigma^{e} \bmod n = H(x) \bmod n$. In the login the signed data is
\[
  x = \texttt{"rsh-auth"} \,\Vert\, K \,\Vert\, \textit{user}.
\]
Because $K$ is fresh for every connection and only the real server can know it,
the signature proves that the client holds the private key \emph{now}, and it
cannot be moved to a connection with a different session key. The prefix
\texttt{"rsh-auth"} keeps this signature from being valid for any other purpose.

\begin{table}[H]
\centering
\footnotesize
% generated by tools/gen_tables.py from data/ -- do not edit
\begin{tabularx}{\textwidth}{@{}l X@{}}
\toprule
Quantity & Value \\
\midrule
$x$ = \texttt{"rsh-auth"} $\Vert K \Vert$ \texttt{"student"} & \texttt{72 73 68 2d 61 75 74 68} \; \texttt{53 51 46 f5 60 34 ab 63 20 13 66 85 7b 35 96 3e} \; \texttt{df f5 68 45 22 dd 29 25 31 81 e6 d2 53 79 7d d0} \; \texttt{73 74 75 64 65 6e 74} (47 bytes) \\
$f = \mathrm{FNV1a}(x)$ & \texttt{56d60f39a9239a06} \\
$H(x) = \mathrm{rotl}_{17}(\mathrm{FNV1a}(x \Vert f)) \oplus f$ & \texttt{34a72278331f3f67} $= 3\,794\,039\,110\,732\,037\,991$ \\
$h = H(x) \bmod n_c$ & $3\,794\,039\,110\,732\,037\,991$ \\
$\sigma = h^{d_c} \bmod n_c$ (client, private) & $7\,793\,721\,554\,923\,979\,395$ \\
$\sigma^{e} \bmod n_c$ (server, public) & $3\,794\,039\,110\,732\,037\,991$ \quad $= h$, accepted \\
\bottomrule
\end{tabularx}

\caption{The login signature of the recorded session. $n_c$, $d_c$ are the client's key (\texttt{id\_rsh}).}
\label{tab:auth}
\end{table}

% ................................................................
\subsection{The stream cipher}
\label{sec:cipher}

A stream cipher turns a key into a long pseudo-random \emph{keystream}
$k_0, k_1, k_2, \ldots$ and encrypts by XOR:
\[
  \text{ciphertext}_j = \text{plaintext}_j \oplus k_j, \qquad
  \text{plaintext}_j = \text{ciphertext}_j \oplus k_j .
\]
Encryption and decryption are the same operation, and the ciphertext is exactly
as long as the plaintext. Both ends compute the same keystream from $K$, so
nothing but $K$ has to be shared. The generator has four 64-bit state words
$s_0, \ldots, s_3$ and is built in three steps (\texttt{StreamCipher::new}).

\begin{enumerate}
\item \textbf{Direction label.} Each direction has its own label,
  \texttt{"rsh client to server"} or \texttt{"rsh server to client"}. Its
  FNV-1a hash seeds \texttt{splitmix64}, which fills the four state words:
  \begin{align*}
    z &\leftarrow z + \texttt{9e3779b97f4a7c15}, &
    z &\leftarrow (z \oplus (z \gg 30)) \cdot \texttt{bf58476d1ce4e5b9}, \\
    z &\leftarrow (z \oplus (z \gg 27)) \cdot \texttt{94d049bb133111eb}, &
    s_i &= z \oplus (z \gg 31),
  \end{align*}
  all modulo $2^{64}$, once for each $s_i$.
\item \textbf{Key.} The 32 key bytes are read as four little-endian words
  $w_0, \ldots, w_3$, and $s_i \leftarrow s_i \oplus w_i$.
\item \textbf{Warm-up.} The first 16 outputs are thrown away, so that the
  first keystream bytes are not a simple function of the mixed-in key.
\end{enumerate}

\noindent Each step of the generator (the \texttt{xoshiro256**} design) outputs
one 64-bit word and updates the state:
\begin{gather*}
  \text{out} = \mathrm{rotl}_{7}(5 s_1) \cdot 9 \bmod 2^{64}, \qquad t = s_1 \ll 17, \\
  s_2 \mathrel{\oplus}= s_0,\quad
  s_3 \mathrel{\oplus}= s_1,\quad
  s_1 \mathrel{\oplus}= s_2,\quad
  s_0 \mathrel{\oplus}= s_3,\quad
  s_2 \mathrel{\oplus}= t,\quad
  s_3 = \mathrm{rotl}_{45}(s_3).
\end{gather*}
The 8 bytes of each output word, lowest byte first, are the next 8 keystream
bytes. \texttt{apply} keeps the unused bytes of the current word between calls,
so splitting a message into pieces gives the same result
(\texttt{stream\_cipher\_round\_trip\_and\_directions}).

\begin{table}[H]
\centering
\scriptsize
% generated by tools/gen_tables.py from data/ -- do not edit
\begin{tabular}{@{}l l l@{}}
\toprule
Step & client $\to$ server & server $\to$ client \\
& label ``\texttt{rsh client to server}'' & label ``\texttt{rsh server to client}'' \\
\midrule
$\mathrm{FNV1a}(\text{label})$ = seed & \texttt{d3a694d4c48f1b05} & \texttt{d48c44aa9e7fb325} \\
\midrule
$s_0 = \mathrm{splitmix64}(\text{seed})$ & \texttt{4f3019c33ecfe8eb} & \texttt{dad3c3c05e48f4ca} \\
$s_1 = \mathrm{splitmix64}(\text{seed})$ & \texttt{1ceb1441c2b7b205} & \texttt{9b55f9973e137179} \\
$s_2 = \mathrm{splitmix64}(\text{seed})$ & \texttt{71a42ea1166f3ef6} & \texttt{25994cc21980dc33} \\
$s_3 = \mathrm{splitmix64}(\text{seed})$ & \texttt{4b66584f85738ef2} & \texttt{48bda24c8b29de32} \\
\midrule
$w_0$ = key bytes 0--7 (little-endian) & \texttt{63ab3460f5465153} & \texttt{63ab3460f5465153} \\
$w_1$ = key bytes 8--15 (little-endian) & \texttt{3e96357b85661320} & \texttt{3e96357b85661320} \\
$w_2$ = key bytes 16--23 (little-endian) & \texttt{2529dd224568f5df} & \texttt{2529dd224568f5df} \\
$w_3$ = key bytes 24--31 (little-endian) & \texttt{d07d7953d2e68131} & \texttt{d07d7953d2e68131} \\
\midrule
$s_0 \oplus w_0$ & \texttt{2c9b2da3cb89b9b8} & \texttt{b978f7a0ab0ea599} \\
$s_1 \oplus w_1$ & \texttt{227d213a47d1a125} & \texttt{a5c3ccecbb756259} \\
$s_2 \oplus w_2$ & \texttt{548df3835307cb29} & \texttt{00b091e05ce829ec} \\
$s_3 \oplus w_3$ & \texttt{9b1b211c57950fc3} & \texttt{98c0db1f59cf5f03} \\
\midrule
$s_0$ after 16 discarded steps & \texttt{9e6e26e8b4957c75} & \texttt{58cc331f6a7ec43d} \\
$s_1$ after 16 discarded steps & \texttt{a6e3e06651e78d2d} & \texttt{0c7065c7616106b3} \\
$s_2$ after 16 discarded steps & \texttt{20c285c7c238fadd} & \texttt{e046ac21ec27af51} \\
$s_3$ after 16 discarded steps & \texttt{43ae7799161bad16} & \texttt{9340e34bf01b1ac4} \\
\midrule
output word 0 & \texttt{0738fe32d9e875a9} & \texttt{e0f2060f0716bc97} \\
output word 1 & \texttt{f769edf40203329c} & \texttt{0f0ef6d275a81be4} \\
output word 2 & \texttt{a355c33eb4673608} & \texttt{2ea23c45d718e280} \\
output word 3 & \texttt{8a87389cc5e19030} & \texttt{e27474c16e09e173} \\
\midrule
keystream bytes 0--7 & \texttt{a9 75 e8 d9 32 fe 38 07} & \texttt{97 bc 16 07 0f 06 f2 e0} \\
keystream bytes 8--15 & \texttt{9c 32 03 02 f4 ed 69 f7} & \texttt{e4 1b a8 75 d2 f6 0e 0f} \\
\bottomrule
\end{tabular}

\caption{Setting up the two ciphers of the recorded session from $K$. The decoder repeated these steps on its own and got the same keystream as \texttt{StreamCipher}.}
\label{tab:ciphersetup}
\end{table}

The two directions use the same $K$ but different labels, so their keystreams
are unrelated (Table~\ref{tab:ciphersetup}, last two rows). This matters: if both
directions used one keystream, XOR-ing a client byte with the server byte at the
same position would cancel the key and leave the XOR of the two plaintexts.

% ................................................................
\subsection{Encrypting each packet}
\label{sec:packets}
\label{sec:rshsession}

\textbf{The rule.} Every message is sent as one frame:
\[
  \underbrace{\text{len}}_{\text{4 bytes, clear}} \;\Vert\;
  \underbrace{\text{encode}(\text{message}) \oplus k_{o}\,k_{o+1} \ldots k_{o+\text{len}-1}}_{\text{body}},
\]
where $o$ is the number of bytes this direction has already encrypted. Each
side keeps one cipher for sending and one for receiving, and each cipher just
continues where the previous message stopped. No counter or nonce travels with
the frame; both ends stay in step because TCP delivers every byte exactly once
and in order. The length stays in clear because the receiver needs it to find
the end of the frame before it can decrypt anything. The first two frames,
\texttt{HostKey} and \texttt{SessionKey}, are sent before the key exists and
are not encrypted.

\textbf{The recorded session.} The client typed \texttt{whoami},
\texttt{echo hello} and \texttt{exit}, with a pause after each. All 16 frames, in
the order they arrived:

\begin{table}[H]
\centering
\footnotesize
% generated by tools/gen_tables.py from data/ -- do not edit
\begin{tabular}{@{}r l r l l l@{}}
\toprule
\# & Dir. & Bytes & Keystream & Message (decrypted) & Body on the wire \\
\midrule
1 & S$\to$C & 21 & -- & HostKey($n_h$, $e$) & \texttt{01 da a9 69 5e 40 \ldots} \\
2 & C$\to$S & 53 & -- & SessionKey(44-byte RSA ciphertext) & \texttt{02 00 00 00 2c 00 \ldots} \\
3 & C$\to$S & 40 & 0--35 & Auth(student, $n_c$, $e$, $\sigma$) & \texttt{aa 75 e8 d9 35 8d \ldots} \\
4 & S$\to$C & 24 & 0--19 & Success(``welcome student'') & \texttt{93 bc 16 07 00 71 \ldots} \\
5 & C$\to$S & 5 & 36--36 & Shell & \texttt{b1} \\
6 & S$\to$C & 22 & 20--37 & Success(``shell started'') & \texttt{41 3c a2 2e 7e 92 \ldots} \\
7 & S$\to$C & 59 & 38--92 & Data(``sh: 0: can't access tty; job co...'') & \texttt{d1 0b 64 3b 49 07 \ldots} \\
8 & C$\to$S & 16 & 37--48 & Data(``whoami\textbackslash{}n'') & \texttt{f1 b0 b1 3c da 40 \ldots} \\
9 & S$\to$C & 11 & 93--99 & Data(``\# '') & \texttt{19 c2 fd 69 8b 8e \ldots} \\
10 & S$\to$C & 14 & 100--109 & Data(``root\textbackslash{}n'') & \texttt{21 cd dc bd d5 95 \ldots} \\
11 & C$\to$S & 20 & 49--64 & Data(``echo hello\textbackslash{}n'') & \texttt{5a 23 07 66 e3 4a \ldots} \\
12 & S$\to$C & 15 & 110--120 & Data(``hello\textbackslash{}n'') & \texttt{92 f3 aa f4 dc 2d \ldots} \\
13 & S$\to$C & 11 & 121--127 & Data(``\# '') & \texttt{90 fc d4 ef 7b 2e \ldots} \\
14 & C$\to$S & 14 & 65--74 & Data(``exit\textbackslash{}n'') & \texttt{e5 58 27 0f 30 36 \ldots} \\
15 & S$\to$C & 13 & 128--136 & Exit(0) & \texttt{38 d7 81 fb 9c 6e \ldots} \\
16 & C$\to$S & 5 & 75--75 & Disconnect & \texttt{60} \\
\bottomrule
\end{tabular}

\caption{Every frame of the recorded session. ``Keystream'' gives the positions of the keystream bytes used for the body.}
\label{tab:timeline}
\end{table}

Frames~1--4 are the handshake of Figure~\ref{fig:handshake}. In frame~5 the
client asks for a shell, and in frame~6 the server confirms that the shell
started. Frames~7--16 are the interactive phase. Frame~7 is a warning that
\texttt{sh} printed to standard error, and frames~9 and~13 are its prompt
\texttt{\#~}. Each typed line is one \texttt{Data} message (frames 8, 11 and 14)
and each piece of output another (10 and 12). After \texttt{exit} the server
sends \texttt{Exit(0)} (frame 15), and the client answers \texttt{Disconnect}
(frame 16). The tables below
show how five of these frames were encrypted.

\textbf{Frame 3: \texttt{Auth}.} The first encrypted frame of the client, so its
body uses keystream bytes 0--35, the first 36 bytes of the client$\to$server
stream in Table~\ref{tab:ciphersetup}.

\begin{table}[H]
\centering
% generated by tools/gen_tables.py from data/ -- do not edit
{\footnotesize length prefix in clear: \texttt{00 00 00 24} $= 36$; keystream bytes 0--35 of the client$\to$server stream\par\smallskip}
{\setlength{\tabcolsep}{3pt}\footnotesize
\begin{tabular}{@{}lcccccccccccc@{}}
\toprule
field & {\scriptsize\itshape tag} & \multicolumn{4}{c}{\scriptsize\itshape len} & \multicolumn{7}{c}{\scriptsize\itshape user} \\
plaintext & \texttt{03} & \texttt{00} & \texttt{00} & \texttt{00} & \texttt{07} & \texttt{73} & \texttt{74} & \texttt{75} & \texttt{64} & \texttt{65} & \texttt{6e} & \texttt{74} \\
as text & \texttt{$\cdot$} & \texttt{$\cdot$} & \texttt{$\cdot$} & \texttt{$\cdot$} & \texttt{$\cdot$} & \texttt{s} & \texttt{t} & \texttt{u} & \texttt{d} & \texttt{e} & \texttt{n} & \texttt{t} \\
$\oplus$ keystream & \texttt{a9} & \texttt{75} & \texttt{e8} & \texttt{d9} & \texttt{32} & \texttt{fe} & \texttt{38} & \texttt{07} & \texttt{9c} & \texttt{32} & \texttt{03} & \texttt{02} \\
= on the wire & \texttt{\textbf{aa}} & \texttt{\textbf{75}} & \texttt{\textbf{e8}} & \texttt{\textbf{d9}} & \texttt{\textbf{35}} & \texttt{\textbf{8d}} & \texttt{\textbf{4c}} & \texttt{\textbf{72}} & \texttt{\textbf{f8}} & \texttt{\textbf{57}} & \texttt{\textbf{6d}} & \texttt{\textbf{76}} \\
\midrule
field & \multicolumn{8}{c}{\scriptsize\itshape $n_c$} & \multicolumn{4}{c}{\scriptsize\itshape $e$} \\
plaintext & \texttt{78} & \texttt{24} & \texttt{e3} & \texttt{af} & \texttt{bf} & \texttt{97} & \texttt{16} & \texttt{6f} & \texttt{00} & \texttt{00} & \texttt{00} & \texttt{00} \\
as text & \texttt{x} & \texttt{\$} & \texttt{$\cdot$} & \texttt{$\cdot$} & \texttt{$\cdot$} & \texttt{$\cdot$} & \texttt{$\cdot$} & \texttt{o} & \texttt{$\cdot$} & \texttt{$\cdot$} & \texttt{$\cdot$} & \texttt{$\cdot$} \\
$\oplus$ keystream & \texttt{f4} & \texttt{ed} & \texttt{69} & \texttt{f7} & \texttt{08} & \texttt{36} & \texttt{67} & \texttt{b4} & \texttt{3e} & \texttt{c3} & \texttt{55} & \texttt{a3} \\
= on the wire & \texttt{\textbf{8c}} & \texttt{\textbf{c9}} & \texttt{\textbf{8a}} & \texttt{\textbf{58}} & \texttt{\textbf{b7}} & \texttt{\textbf{a1}} & \texttt{\textbf{71}} & \texttt{\textbf{db}} & \texttt{\textbf{3e}} & \texttt{\textbf{c3}} & \texttt{\textbf{55}} & \texttt{\textbf{a3}} \\
\midrule
field & \multicolumn{4}{c}{\scriptsize\itshape $e$} & \multicolumn{8}{c}{\scriptsize\itshape $\sigma$} \\
plaintext & \texttt{00} & \texttt{01} & \texttt{00} & \texttt{01} & \texttt{6c} & \texttt{28} & \texttt{dc} & \texttt{76} & \texttt{17} & \texttt{8c} & \texttt{82} & \texttt{83} \\
as text & \texttt{$\cdot$} & \texttt{$\cdot$} & \texttt{$\cdot$} & \texttt{$\cdot$} & \texttt{l} & \texttt{(} & \texttt{$\cdot$} & \texttt{v} & \texttt{$\cdot$} & \texttt{$\cdot$} & \texttt{$\cdot$} & \texttt{$\cdot$} \\
$\oplus$ keystream & \texttt{30} & \texttt{90} & \texttt{e1} & \texttt{c5} & \texttt{9c} & \texttt{38} & \texttt{87} & \texttt{8a} & \texttt{3c} & \texttt{f0} & \texttt{3f} & \texttt{c6} \\
= on the wire & \texttt{\textbf{30}} & \texttt{\textbf{91}} & \texttt{\textbf{e1}} & \texttt{\textbf{c4}} & \texttt{\textbf{f0}} & \texttt{\textbf{10}} & \texttt{\textbf{5b}} & \texttt{\textbf{fc}} & \texttt{\textbf{2b}} & \texttt{\textbf{7c}} & \texttt{\textbf{bd}} & \texttt{\textbf{45}} \\
\bottomrule
\end{tabular}}

\caption{Encryption of the \texttt{Auth} message: user \texttt{student}, the client's public key $(n_c, e)$ and the signature $\sigma$ of Table~\ref{tab:auth}.}
\end{table}

\textbf{Frame 5: \texttt{Shell}.} The body is a single tag byte, \texttt{06}.
The client$\to$server stream is at position 36 after the 36-byte \texttt{Auth}
body, so it uses one keystream byte:

\begin{table}[H]
\centering
% generated by tools/gen_tables.py from data/ -- do not edit
{\footnotesize length prefix in clear: \texttt{00 00 00 01} $= 1$; keystream bytes 36--36 of the client$\to$server stream\par\smallskip}
{\setlength{\tabcolsep}{3pt}\footnotesize
\begin{tabular}{@{}lcccccccccccc@{}}
\toprule
field & {\scriptsize\itshape tag} & & & & & & & & & & & \\
plaintext & \texttt{06} & & & & & & & & & & & \\
as text & \texttt{$\cdot$} & & & & & & & & & & & \\
$\oplus$ keystream & \texttt{b7} & & & & & & & & & & & \\
= on the wire & \texttt{\textbf{b1}} & & & & & & & & & & & \\
\bottomrule
\end{tabular}}

\caption{Encryption of the \texttt{Shell} request.}
\end{table}

\textbf{Frame 8: typing \texttt{whoami}.} The client's terminal sends the line
with its newline. The \texttt{Data} message is the tag \texttt{07}, a 4-byte
length and the 7 bytes:

\begin{table}[H]
\centering
% generated by tools/gen_tables.py from data/ -- do not edit
{\footnotesize length prefix in clear: \texttt{00 00 00 0c} $= 12$; keystream bytes 37--48 of the client$\to$server stream\par\smallskip}
{\setlength{\tabcolsep}{3pt}\footnotesize
\begin{tabular}{@{}lcccccccccccc@{}}
\toprule
field & {\scriptsize\itshape tag} & \multicolumn{4}{c}{\scriptsize\itshape len} & \multicolumn{7}{c}{\scriptsize\itshape text} \\
plaintext & \texttt{07} & \texttt{00} & \texttt{00} & \texttt{00} & \texttt{07} & \texttt{77} & \texttt{68} & \texttt{6f} & \texttt{61} & \texttt{6d} & \texttt{69} & \texttt{0a} \\
as text & \texttt{$\cdot$} & \texttt{$\cdot$} & \texttt{$\cdot$} & \texttt{$\cdot$} & \texttt{$\cdot$} & \texttt{w} & \texttt{h} & \texttt{o} & \texttt{a} & \texttt{m} & \texttt{i} & \texttt{\textbackslash n} \\
$\oplus$ keystream & \texttt{f6} & \texttt{b0} & \texttt{b1} & \texttt{3c} & \texttt{dd} & \texttt{37} & \texttt{84} & \texttt{c9} & \texttt{c9} & \texttt{d1} & \texttt{5b} & \texttt{3f} \\
= on the wire & \texttt{\textbf{f1}} & \texttt{\textbf{b0}} & \texttt{\textbf{b1}} & \texttt{\textbf{3c}} & \texttt{\textbf{da}} & \texttt{\textbf{40}} & \texttt{\textbf{ec}} & \texttt{\textbf{a6}} & \texttt{\textbf{a8}} & \texttt{\textbf{bc}} & \texttt{\textbf{32}} & \texttt{\textbf{35}} \\
\bottomrule
\end{tabular}}

\caption{Encryption of the typed line \texttt{whoami}.}
\label{tab:xorwhoami}
\end{table}

Compare this with the Telnet line in Table~\ref{tab:telnetpw}: there, the
bytes on the wire were the characters themselves; here they are
\texttt{40 ec a6 a8 bc 32 35}.

\textbf{Frame 10: the answer \texttt{root}.} The shell's output travels in the
other direction, with the server$\to$client keystream:

\begin{table}[H]
\centering
% generated by tools/gen_tables.py from data/ -- do not edit
{\footnotesize length prefix in clear: \texttt{00 00 00 0a} $= 10$; keystream bytes 100--109 of the server$\to$client stream\par\smallskip}
{\setlength{\tabcolsep}{3pt}\footnotesize
\begin{tabular}{@{}lcccccccccccc@{}}
\toprule
field & {\scriptsize\itshape tag} & \multicolumn{4}{c}{\scriptsize\itshape len} & \multicolumn{5}{c}{\scriptsize\itshape text} & & \\
plaintext & \texttt{07} & \texttt{00} & \texttt{00} & \texttt{00} & \texttt{05} & \texttt{72} & \texttt{6f} & \texttt{6f} & \texttt{74} & \texttt{0a} & & \\
as text & \texttt{$\cdot$} & \texttt{$\cdot$} & \texttt{$\cdot$} & \texttt{$\cdot$} & \texttt{$\cdot$} & \texttt{r} & \texttt{o} & \texttt{o} & \texttt{t} & \texttt{\textbackslash n} & & \\
$\oplus$ keystream & \texttt{26} & \texttt{cd} & \texttt{dc} & \texttt{bd} & \texttt{d0} & \texttt{e7} & \texttt{1e} & \texttt{56} & \texttt{6f} & \texttt{a9} & & \\
= on the wire & \texttt{\textbf{21}} & \texttt{\textbf{cd}} & \texttt{\textbf{dc}} & \texttt{\textbf{bd}} & \texttt{\textbf{d5}} & \texttt{\textbf{95}} & \texttt{\textbf{71}} & \texttt{\textbf{39}} & \texttt{\textbf{1b}} & \texttt{\textbf{a3}} & & \\
\bottomrule
\end{tabular}}

\caption{Encryption of the shell's output \texttt{root}.}
\end{table}

\textbf{Frame 15: \texttt{Exit(0)}.} The exit code is sent as a 64-bit number:

\begin{table}[H]
\centering
% generated by tools/gen_tables.py from data/ -- do not edit
{\footnotesize length prefix in clear: \texttt{00 00 00 09} $= 9$; keystream bytes 128--136 of the server$\to$client stream\par\smallskip}
{\setlength{\tabcolsep}{3pt}\footnotesize
\begin{tabular}{@{}lcccccccccccc@{}}
\toprule
field & {\scriptsize\itshape tag} & \multicolumn{8}{c}{\scriptsize\itshape code} & & & \\
plaintext & \texttt{09} & \texttt{00} & \texttt{00} & \texttt{00} & \texttt{00} & \texttt{00} & \texttt{00} & \texttt{00} & \texttt{00} & & & \\
as text & \texttt{$\cdot$} & \texttt{$\cdot$} & \texttt{$\cdot$} & \texttt{$\cdot$} & \texttt{$\cdot$} & \texttt{$\cdot$} & \texttt{$\cdot$} & \texttt{$\cdot$} & \texttt{$\cdot$} & & & \\
$\oplus$ keystream & \texttt{31} & \texttt{d7} & \texttt{81} & \texttt{fb} & \texttt{9c} & \texttt{6e} & \texttt{07} & \texttt{90} & \texttt{71} & & & \\
= on the wire & \texttt{\textbf{38}} & \texttt{\textbf{d7}} & \texttt{\textbf{81}} & \texttt{\textbf{fb}} & \texttt{\textbf{9c}} & \texttt{\textbf{6e}} & \texttt{\textbf{07}} & \texttt{\textbf{90}} & \texttt{\textbf{71}} & & & \\
\bottomrule
\end{tabular}}

\caption{Encryption of the shell's exit code.}
\end{table}

\textbf{Decryption.} The receiver reads the 4-byte length, reads that many
bytes, XORs them with the next bytes of its receiving keystream, which are the
same bytes the sender used, and decodes the message. \texttt{Message::decode}
then checks that the tag is known, that every length fits and that no bytes are
left over, so a frame decrypted with the wrong keystream usually fails there
instead of producing a wrong message.

% ................................................................
\subsection{Key files}

Keys are stored as one line of text each, with numbers in hexadecimal. The
files of the recorded session, all in \texttt{\$RSH\_HOME}:

\begin{lstlisting}[style=plain]
host_key          rsh-rsa-private daa9695e409806ab 10001 a7b5ec83dc4684d1   # n, e, d
host_key.pub      rsh-rsa daa9695e409806ab 10001 host                       # n, e, comment
id_rsh.pub        rsh-rsa 7824e3afbf97166f 10001 student@labpc
authorized_keys   rsh-rsa 7824e3afbf97166f 10001 student@labpc              # one client per line
known_hosts       127.0.0.1:2401 rsh-rsa daa9695e409806ab 10001             # host:port, then key
\end{lstlisting}

\noindent Here $\texttt{daa9695e409806ab}_{16} = 15\,756\,240\,624\,863\,938\,219$
is $n$ and $\texttt{a7b5ec83dc4684d1}_{16} = 12\,084\,825\,226\,199\,074\,001$
is $d$ of Table~\ref{tab:hostkey}.

% ................................................................
\subsection{Limits of this design}
\label{sec:cryptolimits}

The scheme is complete in structure but weak in its parts. Each weakness below
was shown by a test (Section~\ref{sec:tests}) or by the recorded session.

\begin{itemize}
\item \textbf{The modulus can be factored.} With $n < 2^{64}$, Pollard's rho
  finds $p$ in about a millisecond (Table~\ref{tab:factor}). From $p$ and $q$
  follow $\varphi$ and $d$, from $d$ the session key in
  Table~\ref{tab:rsablocks}, and from the session key the whole recorded
  session. The decoder used in this section does the same, except that it was
  given $d$ instead of computing it.
\item \textbf{Known plaintext gives keystream away.} Many plaintext bytes can be
  guessed: tags, length fields, prompts. In Table~\ref{tab:xorwhoami} the
  length field is \texttt{00 00 00 07}, so the wire bytes \texttt{b0 b1 3c}
  \emph{are} keystream bytes 38--40. The generator was designed for simulations,
  not for cryptography, and its state can be computed from enough known
  output.
\item \textbf{No integrity.} Flipping a ciphertext bit flips the same plaintext
  bit. Knowing only the position, an attacker turned \texttt{rm -i} into
  \texttt{rm -f} (\texttt{flipped\_ciphertext\_bit\_\allowbreak flips\_plaintext\_bit}).
  Real protocols add a message authentication code to every frame.
\item \textbf{No server randomness.} Only the client chooses $K$, so a recorded
  \texttt{SessionKey} and \texttt{Auth} can be sent again and are accepted
  (\texttt{recorded\_login\_can\_be\_replayed}).
\item \textbf{Textbook RSA.} Blocks are encrypted without random padding, so
  the same 7 bytes under the same key always give the same 8 bytes. With a
  random 32-byte key this leaks nothing here, but it would for guessable data.
\item \textbf{A 64-bit seed.} The session key comes from a generator with 64
  bits of state, so a 32-byte key has at most $2^{64}$ possible values.
\end{itemize}
